%% Composed proof for the active clausal-form theorem; distinct from
%% any unheld original source presentation.
\begin{annexitem}{annex:p2.fol.clausal.quantifier-laws}{thm:p02-fol-quantifier-equivalences}{All twelve quantifier laws and their domain conditions}
\begin{proof}
Fix a nonempty-domain structure and assignment.  For every row with
\(x\notin\FV(\chi)\), coincidence of satisfaction makes the truth
of \(\chi\) independent of updating \(x\).  The following cases
therefore concern that same truth value under every update; this is
where the free-variable side condition enters the argument.

For Q1, failure of \(\forall x\varphi\) means that some domain
element makes \(\varphi\) false, exactly the existential clause
for \(\exists x\neg\varphi\).  For Q2, failure of
\(\exists x\varphi\) means that every domain element makes
\(\varphi\) false, exactly \(\forall x\neg\varphi\).

For Q3, if \(\chi\) is true, both sides reduce to
\(Qx\varphi\), for either quantifier.  If it is false, the left
side is false.  The existential right side is false because no
update makes the conjunction true.  The universal right side is
false because the domain has an element and the conjunction is
false there.  For Q4, if \(\chi\) is false, both sides reduce
to \(Qx\varphi\).  If it is true, the left side is true, the
universal right side is true at every update, and the existential
right side is true because the domain has an element.

For Q5, a true \(\chi\) makes the left implication true and
makes every update of the right matrix true; nonemptiness supplies
an existential witness.  A false \(\chi\) makes the left side
\(\neg\forall x\varphi\) and the right side
\(\exists x\neg\varphi\), already equated in Q1.  For Q6,
a true \(\chi\) makes both sides true; a false \(\chi\)
reduces them to \(\neg\exists x\varphi\) and
\(\forall x\neg\varphi\), equated in Q2.

For Q7, a false \(\chi\) makes both the left implication and
every right-matrix instance true.  A true \(\chi\) reduces
both sides to \(\forall x\varphi\).  For Q8, a false
\(\chi\) makes the left side true and makes each right-matrix
instance true, with a witness supplied by nonemptiness.  A true
\(\chi\) reduces both sides to \(\exists x\varphi\).
These computations prove the implications in both directions for
each displayed law; they do not move an \(x\)-dependent formula
outside a quantifier.

For Q9, every update makes a conjunction true exactly when every
update makes its first conjunct true and every update makes its
second conjunct true.  For Q10, an update making a disjunction true
makes at least one disjunct true; conversely, a witness to either
existential disjunct witnesses their disjunction.  No shared-witness
conjunction assertion is used here.

For Q11 and Q12, coincidence makes \(\varphi\) constant in
the updated variable.  If it is true, every update satisfies it,
and a domain element supplies the existential witness.  If it is
false, no update satisfies it and a domain element refutes the
universal formula.  Thus each quantified formula has exactly the
truth value of \(\varphi\).

The nonempty-domain assumption cannot simply be removed from all
rows.  With an empty domain, \(\forall x\bottom\) is true
whereas \(\bottom\) is false, and
\(\exists x\neg\bottom\) is false whereas
\(\neg\bottom\) is true.  The comparison uses closed formulas
and the usual extension of the satisfaction clauses to empty-domain
sentence semantics.  It does not introduce an empty-domain model
into the standing first-order semantics of the theorem.
\end{proof}
\end{annexitem}

\begin{annexitem}{annex:p2.fol.clausal.matrix}{thm:p02-matrix-clause-equivalence}{Pointwise correctness of matrix conversion}
\begin{proof}
The recursion in Definition~\ref{def:p02-matrix-clause-construction}
is well founded by formula complexity.  A singleton literal family and
the two unit families are finite.  Union and the image of a finite
Cartesian product under union are finite.  Induction therefore proves
finiteness.  Each literal in an output already occurs in the input;
no operation changes a term or introduces a symbol or variable.

Fix a structure and assignment.  For a literal, the one singleton clause
is true exactly when that literal is true.  For \(\bottom\), the
unique clause is empty and false.  For \(\neg\bottom\), there are
no clauses and the family is true.  These establish all base cases.

At a conjunction, truth of every clause in the union is equivalent to
truth of every clause in each family.  By the two induction hypotheses,
this is equivalent to truth of both conjuncts.

At a disjunction, suppose first that \(\alpha\) is true.  Every
\(C\in\operatorname{Cl}(\alpha)\) then contains a true literal,
which remains in every \(C\cup D\); hence every output clause is
true.  If \(\beta\) is true, use the true literal from \(D\)
instead.  Conversely, if both matrices are false, the induction
hypotheses supply a clause \(C\) in the first family with no true
literal and a clause \(D\) in the second with no true literal.
Their union is an output clause with no true literal, so the output
family is false.  Selecting these two witnesses is a finite logical
existence step, not an application of AC.  This reasoning also covers
empty families: a family with no clauses cannot have a false clause,
and the corresponding matrix is true by induction.  Empty clauses
cause no exception because their union follows the same literal test.
The base, conjunction, and disjunction cases exhaust quantifier-free
NNF.  The asserted equivalence follows for every assignment.
\end{proof}
\end{annexitem}

\begin{annexitem}{annex:p2.fol.clausal.composition}{thm:p02-fol-clausal-form}{Sentence-to-clause composition with expansions and reducts}
\begin{proof}
Let the input sentence have signature \(\mathcal L\).
The fully supplied construction in
Theorem~\ref{thm:prenex-normal-form}, using
Lemma~\ref{lem:fol-nnf}, gives a logically equivalent rectified
prenex sentence whose matrix is quantifier-free NNF.  Its prefix
variables are pairwise distinct; because the input is a sentence,
all variables of the matrix are among those prefix variables.
The NNF, rectification, and prenex operations introduce no nonlogical
symbol and use only finite syntactic recursion and a fixed infinite
variable enumeration.  Their satisfaction equivalences hold under
every assignment in the stated nonempty-domain semantics.

Apply the construction of Theorem~\ref{thm:skolem-equisatisfiable}
to this rectified prenex sentence.  Its output has the form
\(\forall u_1\cdots\forall u_p\,\theta\) in a signature
\(\mathcal L^*\), with one fresh function symbol for each
existential prefix position.  There are finitely many such positions;
an initial existential position gives a constant, that is, arity zero.
Each new term uses only preceding universal variables.  Therefore
\(\theta\) is quantifier-free NNF and
\(\FV(\theta)\subseteq\{u_1,\ldots,u_p\}\).
The inherited Skolem proof constructs, under AC, an expansion of
each model of the prenex sentence satisfying this universal sentence.
In the reverse direction it uses the interpreted term values as
witnesses and then takes the \(\mathcal L\)-reduct, without AC.

Set \(S=\operatorname{Cl}(\theta)\).  By
Theorem~\ref{thm:p02-matrix-clause-equivalence}, for every
\(\mathcal L^*\)-structure and assignment, satisfaction of
\(\theta\) is equivalent to truth of every clause of \(S\).
This step stays in the same expanded signature and is logical
equivalence, not another Skolem or Tseitin step.

We now justify the clause closures, rather than merely erase a
written prefix.  If a structure satisfies
\(\forall u_1\cdots\forall u_p\,\theta\), fix a clause of
\(S\) and any assignment to its free variables.  Those variables
are among the \(u_i\).  Extend their values to the remaining
prefix variables by one fixed domain element; the domain is nonempty.
The universal sentence and the pointwise lemma make this clause true
under that extension.  Coincidence of satisfaction on free variables
makes it true under the original clause assignment.  Thus every
clause's universal closure is satisfied.  This finite extension uses
no choice principle.

Conversely, suppose the structure satisfies every universal clause
closure.  Fix any tuple of values for the prefix variables.  Its
restriction to the free variables of each clause makes that clause
true.  The pointwise lemma then makes \(\theta\) true at the
tuple.  Since the tuple was arbitrary, the universal matrix sentence
holds.  When \(p=0\), there is one empty tuple and the same proof
applies.  When \(S\) is empty, all clause closures and all matrix
assignments are true; when it contains the empty clause, both the
clause-model condition and the universal matrix sentence fail in a
nonempty domain.  These conventions therefore preserve all unit cases.

Starting from a model of the input sentence, the same-signature
normal-form equivalences, the AC-dependent Skolem expansion, and the
just-proved matrix/closure equivalence give an \(S\)-model.
Starting from an \(S\)-model, reverse the matrix/closure equivalence,
apply the choice-free Skolem reduct implication, and reverse the
same-signature normal-form equivalences.  The reduct satisfies the
input sentence.  This proves both satisfiability directions, the
finite-expansion claim, and the stated location of the only choice
assumption.  It does not claim same-signature equivalence across
Skolemization or apply to an unspecified closure of an open formula.
\end{proof}
\end{annexitem}
